\documentclass[11pt]{article}
\usepackage[a4paper,margin=25mm]{geometry}
\usepackage{amsmath,amssymb}
\usepackage{enumitem}
\usepackage[hidelinks]{hyperref}
\usepackage{xurl}
\setlist[description]{font=\normalfont\bfseries,style=nextline,
  leftmargin=0pt,labelsep=0pt,labelwidth=0pt,itemsep=8pt,parsep=2pt}
\newcommand{\src}[2]{\textsf{#1}, line~#2}
\newcommand{\where}[1]{\par{\small\textit{Source:} #1.}}
\newcommand{\R}{\mathbb R}
\newcommand{\E}{\mathbb E}
\title{Reading glossary for the \(O(3)\) continuum-limit paper}
\author{}
\date{Source locators checked October 13, 2026}
\begin{document}
\maketitle

This glossary explains the specialized terminology of
\emph{The canonical massive continuum limit of the two-dimensional
\(O(3)\) model} (October 4, 2026), principally the density class and
fluctuation integration in Sections 3--5.  It also includes the nearby
trajectory terminology.  Entries give the place where a term is defined
or its operational meaning is specified.  Some informal terms have no
standalone definition; these are identified below.  Several geometric
definitions are imported from the September 23 \(O(4)\) paper, reference
[28] of the \(O(3)\) paper.  References to that construction concern its
geometry and bookkeeping, not its quaternion-specific spin formulas.

\section*{Source-file key}

All paths are relative to the repository root.  A locator such as
\src{S}{215} means line 215 of the file assigned key S.  Line numbers
refer to the checked source snapshot; nearby section headings and
formulas help locate the passage if the source changes.

The \(O(3)\) source root, abbreviated \textsf{O3}, is
\begin{quote}\small
\path{preprints/The-canonical-massive-continuum-limit-of-the-two-dimensional-O3-model-October-4-2026/build/}
\end{quote}
The imported \(O(4)\) source root, abbreviated \textsf{O4}, is
\begin{quote}\small
\path{preprints/Sharp-mass-bounds-for-the-two-dimensional-O4-model-September-23-2026/build/}
\end{quote}
\begin{description}[itemsep=2pt,style=standard,leftmargin=16mm]
\item[S] O3/\path{sections/setup.tex}: Section 3, retained densities and locality norms.
\item[R] O3/\path{sections/rg.tex}: Section 5, exact renormalization step.
\item[F] O3/\path{sections/free.tex}: Section 4, block observations and free kernels.
\item[Sh] O3/\path{sections/shooting.tex}: initialization and shooting.
\item[T] O3/\path{sections/trajectories.tex}: offsets, admission, and retuning.
\item[G4] O4/\path{supplements/geometry-rg.tex}: Appendix C, local sector factors and complete supports.
\item[R4] O4/\path{sections/renormalization.tex}: imported RG construction and connected reassembly.
\end{description}

\section*{Scales and geometric objects}
\begin{description}
\item[Layer; regular step; inclined step.]
A layer is one lattice scale.  A step integrates from layer \(j\) to
\(j-1\).  A regular step uses square blocks of side \(L\); an inclined
first step uses rotated blocks of side \(L_*=5L\), followed by regular
steps.  Inclined blockings are used to compare orientations.
\where{\src{R}{14}; \src{F}{25} (cells and weighted means).}

\item[Reference scales \(H_j,p_j,t_j,T_j,M\).]
\(H_j=H+(\log L)/(2\pi)\,j\) is the reference coupling;
\(p_j=(\log H_j)^{P_0}\) is the Gaussian cutoff size;
\(t_j=H_j^{-1/2}p_j\) is the small-chord threshold;
\(T_j=H_j^{-49/100}\) is a larger auxiliary chord threshold;
\(M=\mathfrak m_j=\lceil(\log H_j)^2\rceil\) is the mesoscopic
geometric scale.  The actual coupling \(b\) lies near \(H_j\), and
\(g=b^{-1/2}\), so \(t\asymp gp\).
\where{\src{S}{31}, equation with label \texttt{setup:scale-definitions}.}

\item[Chord; true bad edge (bad bond).]
The chord is the Euclidean spin difference
\(r_e=|q_x-q_y|\) for \(e=\langle x,y\rangle\).  A true bad edge
satisfies \(r_e>t_j\); the set of all such edges is \(D(q)\).
This is distinct from the lower threshold at which a smooth seed
profile can be nonzero.
\where{\src{S}{178}; compare \src{R}{150}.}

\item[Anchor; seed anchor; block anchor.]
An anchor is a reference position used to locate a labeled term and
its connecting paths.  A seed anchor is the position assigned to its
edge, observation, or output-edge test; neighborhoods of seeds are
neighborhoods of these positions.  A block anchor embeds a coarse site
in the fine lattice.  The sources specify these roles and equivariant
assignment rules, rather than a universal extra formula selecting an
endpoint of every edge.  A seed anchor is a spatial position, not a
spin value.
\where{\src{S}{215} and \src{S}{578}; \src{R}{137};
\src{G4}{31} (typed seed anchors).}

\item[Neighborhood \(N_r(P)\).]
The closed radius-\(r\) neighborhood of a set \(P\) of seed anchors.
All distances in the sector construction use the same fine-lattice
metric, with coarse positions embedded at their block anchors.
\where{\src{G4}{16} and \src{G4}{31}; \src{R}{137}.}

\item[Profile; default; failure.]
A profile is a smooth cutoff \(G\) between zero and one.  An edge
profile equals one for \(r_e\le t/4\) and zero for \(r_e\ge t/2\).
An observation profile similarly tests block alignment.  In the exact
identity \(1=G+(1-G)\), the retained good factor \(G\) is called a
default; its complement is a failure.  Later expansions also use signed
failures \(G-1\) when writing a default as \(1+(G-1)\).
\where{\src{R}{150}; \src{G4}{67}; \src{R}{584}.}

\item[Seed; primary.]
A chosen complementary input-edge or observation profile is a seed,
also called a selected primary in the imported construction.  Each
true output bad edge is a mandatory seed as well.  An input seed can
already occur at \(r_e\ge t/4\); it need not be a true bad edge
\(r_e>t\).
\where{\src{R}{150}; \src{G4}{77}.}

\item[Sector; primary pattern; empty pattern.]
A sector is one term in the product expansion of every profile plus
its complement.  Its primary pattern records the selected seeds and
their types, subject to the required output seeds.  The empty pattern
has no seeds.  Geometry is held fixed throughout its Gaussian integral;
the pattern is not recomputed from each Gaussian field value.
\where{\src{R}{150}; \src{G4}{58} and \src{G4}{77}.}

\item[Core.]
Join seeds whenever their radius-\(50M\) neighborhoods intersect.
A connected component, together with its assigned factors, is a core.
A core packages a large-gradient region and the nearby terms that
must be estimated together.  It is more than just the set of its seed
sites.
\where{\src{R}{155}; \src{G4}{31}.}

\item[Footprint.]
The radius-\(50M\) neighborhood \(N_{50M}(P_c)\) of a core's seed set
\(P_c\).  Distinct cores have disjoint footprints.  Footprint
compatibility implements maximality of seed components without adding
a distant maximality test to an individual core factor.
\where{\src{R}{164}; \src{G4}{31} and \src{G4}{58}.}

\item[Owner; assigned factor or row.]
The core to which an object is allocated for integration and energy
accounting.  Different objects have different assignment radii:
\(3M\) for frozen spins, \(5M\) for numerical energy rows,
\(8M\) for stationary squares, and \(14M\) for defaults.
The enlarged sets do not give extra copies of the same energy term.
\where{\src{R}{159}; \src{G4}{36} (ownership table).}

\item[Frozen spin; exterior or nonfrozen site.]
Sites within \(3M\) of the seeds are frozen for the Gaussian
completion.  Their spins are held fixed during that conditional
calculation, then integrated with their original Haar measure.
The other sites form the exterior \(E(P)\) and are Gaussianized.
``Frozen'' does not mean deleted from the full integral.
\where{\src{R}{159} and \src{R}{469}; \src{G4}{361}.}

\item[Hole; hole-status test.]
An informal name for the frozen region removed from the Gaussian
index set: \(E(P)=\Lambda\setminus N_{3M}(P)\).  A hole-status test
asks whether a site has been removed, so it determines which indices
belong to a window matrix.  The source uses this terminology through
its geometry and read rules rather than a separate definition.
\where{\src{G4}{312}, \src{G4}{361}, and \src{G4}{396}.}

\item[Halo.]
An enlarged neighborhood attached to a factor's raw read positions
to include all indirect dependencies.  In the imported connected
construction each noncore read position receives a \(20M\) halo,
large enough to contain its windows and hole-status queries.
It is a support margin, not an additional interaction or field.
\where{\src{R4}{892}; \src{G4}{350}.}

\item[Inverse window \(\Omega_x\).]
The radius-\(R\) neighborhood of a nonfrozen center \(x\), with
frozen sites removed:
\(\Omega_x=E(P)\cap B_R(x)\).  The radius \(R\) is a small
fixed multiple of \(M\), chosen so that local compositions fit within
the prescribed margins.  Matrix inversion takes place on this set.
\where{\src{R}{347}; explicit definition in \src{G4}{369};
radius convention in \src{G4}{16}.}
\end{description}

\section*{Supports, masks, and covering constraints}
\begin{description}
\item[Label.]
A formula packaged with the data needed to use and sum it: an anchor
when assigned, a complete support, and a connected record.  Covering
labels carry an inventory as well.  Two copies of a similar formula
can have different labels because their dependency records differ.
\where{\src{S}{215}; covering labels in \src{S}{515}.}

\item[Read; raw read; read set.]
A spin, Gaussian coordinate, mask, status, or other datum queried in
evaluating a factor.  Raw positions are the initial arguments before
adding windows, tests, and protective margins.  A Gaussian read set
\(I_c\) lists the Gaussian coordinates used by a core.  These are
operational terms whose full dependency rules appear in the read table.
\where{\src{R}{474}; \src{G4}{308} (complete-read table).}

\item[Complete support.]
Every position needed to determine a factor, including indirect
dependencies: spin arguments, masks, eligibility and absence tests,
statuses, inverse windows, and joining paths.  It can be much larger
than the sites visibly occurring in the factor's displayed formula.
Disjoint complete supports are what justify exact factorization.
\where{\src{S}{215}; \src{R}{169}; \src{G4}{378}.}

\item[Record; load.]
A record is a connected collection of radius-\(M\) cubes and joining
paths covering the complete support.  Its integer load \(s\ge1\)
controls support size and path length:
\(|\operatorname{supp}|\le CM^2s\) and
\(\operatorname{length}\le CMs\).  Repeated paths are retained, and
unused load may be kept under blocking.  Thus load is not simply the
number of sites.  The core seed count \(s_c\) in the reserve estimate
is specified separately as a number of seeds.
\where{\src{S}{220}; projection convention in \src{S}{235};
core seed count in \src{R}{474}.}

\item[Bulk formula; lifted record.]
A bulk formula is defined on the infinite lattice and folded onto a
torus.  A lifted record retains displacements before they are reduced
modulo periods, allowing a calculation to be compared with the infinite
lattice even when its displayed arguments coincide on the torus.
\where{\src{S}{17} and \src{S}{248}.}

\item[Short; winding.]
A record is short when \(s<c_*\mathcal L_j/M\), where
\(\mathcal L_j\) is the shortest period.  It then has the required
injective lift and agrees with its bulk construction.  Other records
are winding.  A scalar term can still be winding if its construction
used a long or period-dependent record.
\where{\src{S}{248}.}

\item[Kinetic mask.]
The indicator \(m_e\) enabling a nonlocal kinetic row only if
\(r_f\le t\exp(d(e,f)/R_*)\) for every edge \(f\).
Since chords are at most two, this test queries only a finite
neighborhood.  It can be discontinuous as the spins vary.
\where{\src{S}{184}, label \texttt{eq:setup-kinetic-mask}.}

\item[Witness; failed-mask witness.]
An edge whose chord violates a kinetic mask's defining inequality.
Such an edge explains why the row is switched off and supplies an
energy reserve against which its loss is charged.  Every witness
used by the argument is retained in the support record.
\where{mask definition in \src{S}{194}; charging rule in \src{G4}{240};
initialization in \src{Sh}{16}.}

\item[Graph mask; open graph domain.]
For a connected graph \(X\), the mask \(1_X\) requires
\(r_e\le t_j\) on every edge of \(X\).  The larger open domain
\(\mathcal U_{X,j}\) requires \(r_e<4t_j\) there and is where the
regular function's derivatives are controlled.  The masked function
is zero off its mask.
\where{\src{S}{267}, label \texttt{eq:setup-graph-domain}.}

\item[Mask strengthening.]
Replace a graph mask by a stricter one, typically on a larger common
graph needed for symmetry.  The difference is retained as a covering
correction, so the density does not change.  A nonzero correction
must have a bad-edge witness in the strengthened graph.
\where{\src{S}{578}; reassembly in \src{R4}{928}.}

\item[Inventory; mandatory covering sum (covering gas).]
A covering label lists a nonempty inventory \(J_\ell\) of true bad
edges, and its weight vanishes unless all those edges are bad.
The covering sum \(\Xi\) runs over families with disjoint complete
supports whose inventories partition \(D(q)\) exactly.  Every bad
edge must be covered once.  Individual weights need not be positive.
\where{\src{S}{513}, label \texttt{eq:setup-covering-sum}.}

\item[Output flag; output inventory; no-flag extension.]
A flag records a true bad edge in the output field that must be
covered.  A factor with output inventory carries such required edges.
For an empty-inventory regular calculation, output statuses are fixed
to ``no flag,'' clipping is removed, and good masks are set to one
to define a smooth extension.  Discontinuous statuses are not
differentiated across their jumps.
\where{\src{R}{155} and \src{R}{239}; \src{S}{515}.}

\item[Compatibility; status query; eligibility test.]
These specify whether a labeled object or a family of objects is
allowed.  For covering families, compatibility includes disjoint
complete supports and exact inventory coverage; for primary cores it
includes footprint compatibility.  A status query may ask whether
another object is present or absent.  Such decisions are part of the
complete support, even if they do not appear as a spin argument.
\where{\src{S}{215} and \src{S}{525}; \src{R}{167}; \src{G4}{58}.}
\end{description}

\section*{Gaussian completion and estimates}
\begin{description}
\item[Tag; fallback branch.]
Independently in each block, a tag \(\tau_Y\) selects a constituent
site with the block weights.  If the weighted spin mean has length
at least \(1/2\), the observation uses its unit normalization;
otherwise it uses the tagged spin.  This fallback defines a unit
observation direction everywhere.
\where{\src{S}{83}, labels \texttt{setup:unit-mean} and
\texttt{setup:observation}.}

\item[Energy row; stack; ledger.]
A row is an individual component in the decomposition of the energy.
A stack collects the direct gradient \(\sqrt{c_0}\,dq\) and the
additional gradient \(\ell_s dq\).  ``Ledger'' names the exact
identity expressing input energy plus observation energy minus output
energy as the rows \(\mathcal F,\mathcal B,\mathcal C,\mathcal U,
\mathcal N,\mathcal Z\).  These are accounting terms for that displayed
identity, not additional probabilistic objects.
\where{\src{R}{175}, especially the row identity beginning at line 198.}

\item[Selected row; selected assigned row.]
The \(\mathcal F,\mathcal B\) rows farther than \(2.5M\) from
all seeds are selected for affine Gaussian completion.  Every row
incident on a nonfrozen column is selected.  Selection and core
assignment are different tests, so a row can be selected and also
assigned to a core.
\where{\src{R}{223}; assignment radii in \src{R}{159}.}

\item[Principal logarithm; relative tangent coordinate.]
The principal logarithm \(u_x=\operatorname{Log}_{T_x}q_x\) is the
inverse spherical exponential with \(|u_x|<\pi\).  In contrast,
\(y_o(x)=\operatorname{proj}_{q_o^\perp}q_x\) is a globally defined
tangent projection used for polynomial slots.  They agree to first
order at alignment but are different coordinates.
\where{\src{S}{285}; \src{R}{252}.}

\item[Protector; alignment profile; global chart profile.]
Smooth cutoffs attached to factors to keep their coordinate evaluations
inside valid charts.  Alignment profiles additionally enforce the
small-angle region.  The global chart profile removes nonprincipal
exponential preimages when the integral is extended to ambient space.
For an optional exponential, the prescription is
\(1+\chi_{\rm att}(e^V-1)\), with the protector attached to the
selected correction.
\where{\src{R}{230} and \src{R}{257}; detailed preparation in \src{G4}{116}.}

\item[Ambient completion; color.]
Adjoin a normalized scalar Gaussian \(v_x\) and combine it with
the two-dimensional tangent coordinate:
\(\xi_x=u_x+T_xv_x\in\R^3\).  A color means an internal component
of this ambient vector, not a lattice direction.  There are three
Gaussian colors but only two physical spin coordinates; the added
scalar normalization cancels one color in the leading determinant.
\where{\src{R}{266} (Lemma 5.2); \src{R}{386}.}

\item[Affine row; precision \(A\); window matrix.]
The affine reference row is \(D_i\xi+F_i\), with the frozen
coordinates absorbed into its constant part.  The scalar precision
is \(A=D^*D\), acting on every ambient component.  The window
matrix \(A_{\Omega_x}\) is its principal submatrix, inverted after
restriction to the window.
\where{\src{R}{266}; \src{R}{347}; explicit principal restriction
in \src{G4}{361}.}

\item[Resolvent parameter \(a_0\).]
The nonnegative scalar in \((A_{\Omega_x}+a_0I)^{-1}\).  It is
integrated over to represent logarithmic determinants.  It is
distinct from the observation precision \(a=1\), the coupling
\(b\), and a lattice scale.
\where{\src{R}{347} and \src{R}{394}.}

\item[Short covariance \(C_s\); completed mean; stationary residual.]
The local window inverses form \(\Pi_0\).  Set
\(C_s=(\Pi_0+\Pi_0^*)/2\),
\(\mu=-\Pi_0^*D^*F\), and \(r_s=D\mu+F\).
Then \(\xi=\mu+gZ\) with \(Z\sim N(0,C_s\otimes I_3)\).
\(C_s\) has finite range, but differs from the full inverse;
the mean need not solve the full stationary equation exactly.
The square of a residual row is called a stationary square.
\where{\src{R}{360}; \src{R}{802} (canonical vertices).}

\item[Prediction of a spin.]
The spin obtained from the completed mean with the fluctuation set
to zero:
\[
q_x^0=\operatorname{Exp}_{T_x}
(\operatorname{proj}_{T_x^\perp}\mu_x).
\]
Prediction estimates control its gradient and its size relative to
the block reference.
\where{\src{R}{403}, Lemma 5.3.}

\item[Residual chain; residual hop; ratio letter.]
Terms in the inverse, determinant, and Gaussian-density-ratio
corrections caused by window localization.  A hop is a successive
kernel factor along their recorded paths.  Residual chains contain
an exponentially small residual factor, but may be arbitrarily long.
A ratio letter is a prepared factor from the density-ratio exponent;
its Gaussian endpoints and every intermediate window are recorded.
\where{\src{R}{373} and \src{R}{575}.}

\item[Core factor \(B_c\).]
The product belonging to one core after each selected assigned
affine square has been subtracted exactly once.  It includes the
assigned energies, stationary squares, profiles, inventories,
protectors, and normalization compensations.
\where{\src{R}{474}; explicit imported formula in \src{G4}{160},
label \texttt{app:explicit-core}.}

\item[Reserve; charging; normalization cost.]
A reserve is an exponential suppression retained to pay later
combinatorial or analytic losses.  A core starts with a reserve
of order \(e^{-c_Lp^2s_c}\).  Charging assigns a loss to a specified
positive energy contribution, without spending it twice.  Determinants
and scalar normalizations also incur costs paid by this reserve.
\where{\src{R}{486}; detailed charging argument in \src{G4}{230}.}

\item[Prepared factor; letter; opening an exponential.]
A factor after its cutoffs, supports, and majorants have been fixed.
Opening \(e^V\) means the exact identity
\(e^V=1+(e^V-1)\); the selected nontrivial term is a letter.
The preparation specifies which protectors stay attached, so opening
does not create an uncontrolled contribution outside a chart.
\where{\src{R}{230}; eight classes of factors in \src{R}{553}.}

\item[Ordinary letter; tail letter; exceptional factor.]
An ordinary row letter is restricted to Gaussian reads of size at
most \(2p\) and has a positive power of \(g\) in its majorant.
A tail letter requires a read of size at least \(p\) and is paid
by Gaussian rarity.  Exceptional factors also include sufficiently
small truncation and residual corrections; they need not all be
large-field events.
\where{\src{R}{498}, \src{R}{534}, and \src{R}{553};
summed bounds in \src{R}{599}.}

\item[Remote term.]
An old canonical polynomial or regular error whose whole projected
graph is farther than \(14M\) from all seeds.  Its mask can be
replaced by the specified smooth profile for the regular calculation.
A nonremote term meets a core and retains its actual mask and protectors.
\where{\src{R}{563}.}

\item[Majorant; envelope; site-hit bound.]
A majorant is a nonnegative quantity bounding a factor or its
integrated products.  An envelope is often a pointwise bound that
still depends on Gaussian variables.  A site-hit bound sums weighted
majorants of all labels whose supports contain a specified site,
uniformly in that site.  This pays for the number of possible supports.
\where{\src{R}{486} and \src{R}{599}; explicit tree/site-hit
condition in \src{R4}{869}; covering version in \src{S}{540}.}

\item[Marked discrepancy.]
A single chosen difference between two versions of a factor, coming
from changed coupling, history, or retained data.  Joint estimates
must preserve the small size of that difference while bounding all
other factors.  ``Marked'' identifies where the comparison is inserted.
\where{\src{R}{599} (Proposition 5.5); comparison variables in \src{R}{99}.}

\item[Gaussian guard.]
An even smooth cutoff on Gaussian coordinate reads, used in Taylor
and contraction estimates.  Its complement is controlled by Gaussian
tails.  Evenness allows odd fluctuation terms to integrate to zero.
\where{\src{R}{983}.}
\end{description}

\section*{Reassembly, normalization, and trajectories}
\begin{description}
\item[Integrated polymer; connected reassembly.]
Group a connected collection of selected prepared factors, using
intersection of complete supports, and integrate it jointly.  The
result is an integrated polymer retaining the original compatibility,
coverage, and status checks.  Reassembly recovers the exact density
from compatible families of these objects.
\where{\src{R}{770}; explicit construction in \src{R4}{903}.}

\item[Output-empty background; seeded polymer.]
An output-empty polymer carries no mandatory output bad-edge inventory
and belongs to the regular background.  A seeded polymer carries a
nonempty output inventory.  The Mayer logarithm is applied to the
background; seeded components keep their covering weights.
\where{\src{R}{770}; \src{R4}{903}.}

\item[Canonical slots; canonical extraction.]
The designated low-degree polynomial coefficient lists
\((bP_4,I_2,bP_5,I_3,bP_6,I_4,b^{-1}J_2)\), with
\(P_5=I_3=0\) for \(O(3)\).  Subscripts denote tangent-field
degree, not order in \(g\).  Extraction assigns the connected Wick
Taylor coefficients to these slots, the kinetic coupling, and the
bulk scalar, leaving a controlled remainder.
\where{\src{S}{316}; \src{R}{786} (canonical extraction).}

\item[Affine coefficient (affine Hessian); compensation; quadratic packet.]
The affine coefficient measures the second variation on a slowly
varying affine tangent field.  Compensation subtracts the corresponding
nearest-neighbor quadratic term.  A quadratic packet contains both
the original averaged monomial and this subtraction:
\(Q_{r,s}=M_{r,s}-(r\cdot s)S_o\).
The packet is retained as a unit so its cancellation is visible.
\where{\src{S}{350} and \src{S}{475}.}

\item[Regular error; normalization.]
A regular-error label is \(f_X\in C^8(\mathcal U_{X,j})\) with its
graph and record.  Bulk and short errors have zero value at aligned
spins and zero affine Hessian there after the scalar and kinetic terms
are extracted.  Winding errors have different allowances.
\where{\src{S}{417} and \src{S}{493}.}

\item[Retained density; cap; admitted representation.]
The density class consists of a bulk scalar, masked kinetic energy,
canonical polynomial slots, regular errors, and the covering sum
\(\Xi\).  Caps are the prescribed bounds on these coefficient and
function norms.  A representation is admitted when the support,
symmetry, period, and norm conditions hold and
\(b\in[H_j/2,2H_j]\).
\where{\src{S}{162}; density formula in \src{S}{588};
admission in \src{S}{602}.}

\item[Free history.]
The record of scalar block observations already performed, determining
the free precision and interpolation kernels.  It is not the current
layer index.  For regular observations it can be specified by length;
an inclined first observation also requires its geometry to be recorded.
\where{\src{S}{53} and \src{S}{117}.}

\item[Endpoint density; terminal coupling; shooting.]
An endpoint density is the layer-zero retained representation with
its bulk scalar removed.  The terminal coupling is its \(b_0\).
Shooting chooses a microscopic \(b_N=\beta_N\) so that all steps
are admitted and reach a prescribed terminal value, for example
\(b_0=H\).
\where{\src{S}{644}; \src{Sh}{144}, Proposition with label
\texttt{prop:shooting}.}

\item[Bare offset.]
The difference \(s=\beta-\vartheta_N\) between a depth-\(N\)
microscopic coupling and its reference trajectory.  Here
\(\vartheta_j=H_j-(\kappa_\infty/\gamma)\log(H_j/H)\).
This \(s\) is unrelated to the load of a support record.
\where{\src{T}{326}, especially the definition at line 336.}
\end{description}

\section*{A compact guide to the geometry}
The nested neighborhoods have different purposes.  A seed set \(P_c\)
generates frozen sites within \(3M\), ownership regions out to
\(5M,8M,14M\), and a \(50M\) footprint for compatibility.
A hole is the removed frozen part of the Gaussian index set; a halo
is an enlargement of a factor's reads; a core is the whole connected
seed component with its assigned factors.  None of these should be
identified with the complete support of an arbitrary long residual chain.

Two further distinctions matter when using the Gaussian formulas:
the pattern is fixed during a Gaussian integration, although it is
summed over in the full density; and disjoint \emph{complete} supports
give marginal factorization, while arbitrary overlapping conditional
Gaussian laws are not claimed to be independent
(\src{R}{646}; \src{G4}{378}).

\end{document}
